\textbf{COLUNI 2022 -- Questão 31}\par\medskip Leia a notícia a seguir.

\par

\textbf{Um forte terremoto de magnitude 7,2 atinge o Haiti.}

\par

Um terremoto de magnitude 7,2 atingiu o Haiti no dia 14 de agosto de 2021, de acordo com o Serviço Geológico dos Estados Unidos (USGS), e deixou ao menos 300 mortos, cerca de 2.000 feridos e centenas de desaparecidos segundo as autoridades da Defesa Civil do país. O epicentro foi localizado a 12 quilômetros a nordeste de Saint-Louis du Sud, no sul do país. O terremoto ocorreu às 8h29min, hora local (12,29 GMT) e teve uma profundidade de 10 quilômetros, de acordo com o relatório do USGS. [...] No entanto, de acordo com os primeiros dados, na capital não houve danos comparáveis aos do terremoto de 2010. O tremor foi sentido em todo o país e a área afetada pelo terremoto, a sudoeste da ilha, é a mesma que foi atingida pelo furacão Matthew, que deixou centenas de mortos no Haiti em 2016. [...]

\par

Disponível em: https://brasil.elpais.com/internacional/2021-08-14/um-forte-terremoto-de-magnitude-72-atinge-o-haiti.html. Acesso em: 16 ago. 2021 (Adaptado).

\par

Com base nas informações fornecidas pelo texto e nos conhecimentos geográficos, analise as afirmativas a seguir:

\par

I. O Haiti é um país localizado na área de contato de duas placas tectônicas, a placa do Caribe e a placa norte-americana, fato que explica a frequência de terremotos que acontecem nessa região da América Central.\par
II. De população indígena, o Haiti sofre frequentemente com crises humanitárias e políticas, em função dos constantes eventos naturais, resultando numa extrema pobreza, o que atribui ao país o pior IDH do continente americano.\par
III. Com pouco mais de 11 milhões de habitantes e localizado em uma ilha que faz divisa com a República Dominicana, o Haiti ainda está se recuperando de um terremoto de magnitude 7 na escala Richter que ocorreu perto da capital, em 2010, deixando 300.000 mortos.

\par

Considerando as informações apresentadas, está CORRETO o que se afirma em:\begin{enumerate}
\item I e II, apenas.
\item I e III, apenas.
\item II e III, apenas.
\item I, II e III.\end{enumerate}
